\documentclass[a4paper, 12pt, dvipdfmx]{extarticle}

\usepackage{animate}
\usepackage{amsmath, amssymb, amsthm, mathrsfs, amsfonts, dsfont}
\usepackage{ascmac}
\usepackage{bbm}
\usepackage{bm}
\usepackage{breakcites}
\usepackage{calc}
\usepackage[style=base]{caption}
\usepackage{enumerate}
\usepackage[T1]{fontenc}
\usepackage{ifthen}
\usepackage{mathtools}
\usepackage{makecell}
\usepackage{mlmodern}
\usepackage{newtxtext}
\usepackage{optidef}
\usepackage[deluxe]{otf}
\usepackage{physics}
\usepackage{pifont}
\usepackage{setspace}
\usepackage{stfloats}
\usepackage{subcaption}
\usepackage{svg}
\usepackage{tikz}
\usepackage{xparse}
\usepackage[all]{xy}

\usepackage{float}
\usepackage{color}
\usepackage{graphicx}
\usepackage[margin=20truemm]{geometry}
\usepackage[colorlinks=true, allcolors=blue]{hyperref}

\usepackage{listings}

% === Commands ===

\definecolor{cA}{HTML}{0072BD}
\definecolor{cB}{HTML}{EDB120}
\definecolor{cC}{HTML}{77AC30}
\definecolor{cD}{HTML}{D95319}
\definecolor{cE}{HTML}{7E2F8E}
\newcommand{\cAText}[1]{\textcolor{cA}{#1}}
\newcommand{\cBText}[1]{\textcolor{cB}{#1}}
\newcommand{\cCText}[1]{\textcolor{cC}{#1}}
\newcommand{\cDText}[1]{\textcolor{cD}{#1}}
\newcommand{\cEText}[1]{\textcolor{cE}{#1}}
\newcommand{\red}[1]{\textcolor{red}{#1}}
\newcommand{\blue}[1]{\textcolor{blue}{#1}}
\newcommand{\green}[1]{\textcolor{green}{#1}}
\newcommand{\gray}[1]{\textcolor{gray}{#1}}
\newcommand{\black}[1]{\textcolor{black}{#1}}

\newcommand{\st}{\text{ s.t. }}
\newcommand{\Img}[1]{\mathrm{Im}\qty(#1)}
\newcommand{\Ker}[1]{\mathrm{Ker}\qty(#1)}
\newcommand{\Supp}[1]{\mathrm{supp}\qty(#1)}
\newcommand{\Rank}[1]{\mathrm{rank}\qty(#1)}
\newcommand{\floor}[1]{\left\lfloor #1 \right\rfloor}
\newcommand{\ceil}[1]{\left\lceil #1 \right\rceil}
% C++ (https://tex.stackexchange.com/questions/4302/prettiest-way-to-typeset-c-cplusplus)
\newcommand{\Cpp}{C\nolinebreak[4]\hspace{-.05em}\raisebox{.4ex}{\relsize{-3}{\textbf{++}}}}
% https://tex.stackexchange.com/questions/28836/typesetting-the-define-equals-symbol
\newcommand{\defeq}{\coloneqq}
\newcommand{\eqdef}{\eqqcolon}
% https://tex.stackexchange.com/questions/5502/how-to-get-a-mid-binary-relation-that-grows
\newcommand{\relmiddle}[1]{\mathrel{}\middle#1\mathrel{}}
\newcommand{\cmark}{\cCText{\ding{51}}} % check mark
\newcommand{\xmark}{\red{\ding{55}}} % cross mark

\DeclareMathOperator{\Proj}{Proj}
\DeclareMathOperator{\Exp}{Exp}
\DeclareMathOperator{\Hess}{Hess}
\DeclareMathOperator{\Retr}{Retr}
\DeclareMathOperator{\Span}{span}
\DeclareMathOperator{\myGrad}{grad}
\renewcommand{\grad}{\myGrad}

% https://tex.stackexchange.com/questions/564216/newcommand-for-each-letter
\ExplSyntaxOn
\NewDocumentCommand{\definealphabet}{mmmm}{
\int_step_inline:nnn{`#3}{`#4}{
\cs_new_protected:cpx{#1 \char_generate:nn{##1}{11}}{
\exp_not:N #2{\char_generate:nn{##1}{11}}}}}
\ExplSyntaxOff

\definealphabet{bb}{\mathbb}{A}{Z}
\definealphabet{rm}{\mathrm}{A}{Z}
\definealphabet{cal}{\mathcal}{A}{Z}
% \definealphabet{scr}{\mathscr}{A}{Z}
\definealphabet{frak}{\mathfrak}{a}{z}
\definealphabet{frak}{\mathfrak}{A}{Z}

% === Settings ===

% https://qiita.com/rityo_masu/items/efd44bc8f9229e014237
\allowdisplaybreaks[4]

\usetikzlibrary{
  3d,
  fit,
  calc,
  math,
  matrix,
  patterns,
  backgrounds,
  arrows.meta,
  decorations.pathmorphing,
}

\begin{document}

\title{Supplemental Material for 03/10\\Book Reading Seminar}
\author{Hiroki Hamaguchi}
\date{2026/03/10}
\maketitle

\noindent%
Chernoff bound: \href{https://en.wikipedia.org/wiki/Chernoff_bound#Applications}{Wikipedia}.

\begin{quote}
    A simple and common use of Chernoff bounds is for ``boosting'' of randomized algorithms. If one has an algorithm that outputs a guess that is the desired answer with probability $p > 1/2$, then one can get a higher success rate by running the algorithm $n = \log(1/\delta) 2p/(p - 1/2)^2$ times and outputting a guess that is output by more than $n/2$ runs of the algorithm. (There cannot be more than one such guess.) Assuming that these algorithm runs are independent, the probability that more than $n/2$ of the guesses is correct is equal to the probability that the sum of independent Bernoulli random variables $X_k$ that are $1$ with probability $p$ is more than $n/2$. This can be shown to be at least $1-\delta$ via the multiplicative Chernoff bound:

    \[
        \Pr\left[X > \frac{n}{2}\right] \ge 1 - e^{-n(p - 1/2)^2/(2p)} \geq 1-\delta
    \]
\end{quote}

\vskip\baselineskip\noindent%
Hit and Run Sampling: Markov Chains Monte Carlo (MCMC) の一種. 恐らく珍しいタイプ.

\vskip\baselineskip\noindent%
Simulated annealing in convex bodies and an
$\mathcal{O}^*(n^4)$ volume algorithm:\\
\url{https://www.sciencedirect.com/science/article/pii/S0022000005000966/pdf}

\vskip\baselineskip\noindent%
Volume algorithmが間違っている可能性がある + LaTeXが正しく記述されていないので、\\
こちらの論文解説を後半は行います。

\end{document}
